\section{List of the Protocols}

\subsection{Main Protocol: Asynchronous Byzantine Agreement (ABA)} (taken from \cite{Abraham2022Efficient})
Processes start with a bit $\{0,1\}$ and may return a bit.
a program implements ABA if it satisfies three properties: Validity, Agreement, Almost-Sure Termination.
\begin{itemize}
    \item Validity: if a correct process returns $b$ then a correct process had $1-b$ as input.
    \item Agreement: if two correct processes return, then their output are equal.
    \item Termination: if $n-f$ correct processes take part to the protocol then, with probability 1, all correct processes will eventually return.
\end{itemize}

\begin{protocol}\label{prot:SpecABA}
    \lean{ABA.Spec}
    %\leanok
    \begin{algorithm}
        \caption{Ideal Implementation of ABA}
        \begin{algorithmic}
            \State \textbf{Input:} $b \in \{0,1\}$
            \medskip
            \State \textbf{Code for every process:}
            \State send $\langle input, b\rangle$ to the third party 
            \State \textbf{upon} receiving $\langle output, b\rangle$ from third party 
            \State \quad return $b$
            \medskip
            \State \textbf{Code for third party:}
            \State $bind,call \gets \bot,(\bot)^{Id}$
            \State \textbf{upon} receiving $\langle input, b\rangle$ from a process $id$ and $bind = \bot$:
            \State \quad $call(id) \gets b$
            \State \textbf{upon} $|F| + \{id \in Id \setminus F \mid call(id) = b\} \ge n-f$:
            \State \quad $bind,call \gets b,(\bot)^{Id}$
            \State \quad \textbf{if} $\{id \in Id \setminus F \mid call(id) = 1-b\} = 0$:
            \State \quad\quad send $\langle output, b\rangle$ to all
            \State \textbf{upon} $call(id) = \bot$ and $bind \ne \bot$:
            \State \quad $call \gets bind \textbf{ or } \mu$ \Comment{$\mu$ can be any non-Dirac probability distriution over $\{0,1\}$}
        \end{algorithmic}
    \end{algorithm}
\end{protocol}

\begin{protocol}\label{prot:ImplABA}
    \lean{ABA.Impl}
    %\leanok
    \begin{algorithm}
        \caption{Implementation of ABA taken from \cite{Abraham2022Efficient}}
        \begin{algorithmic}
            \State \textbf{Input:} $b \in \{0,1\}$
            \medskip
            \State $r \gets 1$
            \State $decided \gets \bot$
            \Loop
                \State $(b, g) \gets \textsc{GBCA}_r(b)$ \Comment{$b \in \{0,1,\bot\}$, $g \in \{A, B, C\}$}
                \State $c \gets \textsc{WCC}_r$ \Comment{common coin for round $r$}
                \If{$b = \bot$}
                    \State $b \gets c$
                \ElsIf{$g = A$}
                    \State $decided \gets b$ \Comment{ keep participating to help others}
                \EndIf
                \State $r \gets r + 1$
            \EndLoop
            \State \Comment{TODO: verify termination rule}
        \end{algorithmic}
    \end{algorithm}
\end{protocol}

\begin{definition}\label{def:ABAConcrete}
    \lean{ABA.Concrete}
    %\leanok
    \uses{prot:ImplABA, prot:ImplGBCA, def:WCCConcrete}
    $\mathsf{ABA.Concrete} := \mathsf{ABA.Impl}[\mathsf{GBCA.Impl}, \mathsf{WCC.Impl}]$, the instantiation of $\mathsf{ABA.Impl}$ with concrete implementations of the sub-protocols.
\end{definition}

\begin{definition}\label{def:ABAHybrid}
    \lean{ABA.Hybrid}
    %\leanok
    \uses{prot:ImplABA, prot:SpecGBCA, prot:SpecWCC}
    $\mathsf{ABA.Hybrid} := \mathsf{ABA.Impl}[\mathsf{GBCA.Spec}, \mathsf{WCC.Spec}]$, the instantiation of $\mathsf{ABA.Impl}$ with ideal specifications of the sub-protocols.
\end{definition}

\subsection{ABA Sub-Protocols}

\paragraph{Graded Binding Crusader Agreement (GBCA)} (taken from \cite{Abraham2022Efficient,Attiya2023Multi})
Processes start with a binary input $\{0,1\}$ and may return a value in $\{(0,A), (0,B), (\bot,C), (1,B), (1,A)\}$.
A program implements GBCA if it satisfies four properties: Validity, Graded Agreement, Binding, Termination.
\begin{itemize}
    \item Validity: if a correct process does not return $(b,A)$ then a correct process had $1-b$ as input.
    \item Graded Agreement: if two correct processes return $(b,X)$ and $(b',X')$ then $b = 0 \implies b' \ne 1$ and $X = A \implies X' \ne \bot$.
    \item Binding: when a correct process returns then a bound value $b$ is defined and, in all extensions of the execution, correct processes do not return $(1-b,B)$ or $(1-b,A)$.
    \item Termination: if $n-f$ correct processes take part to the protocol then all correct processes will eventually return.
\end{itemize}

\begin{protocol}\label{prot:SpecGBCA}
    \lean{GBCA.Spec}
    %\leanok
    Blablabla
    \begin{algorithm}
        \caption{Ideal Implementation of GBCA}
        \begin{algorithmic}
            \State \textbf{Input:} $b \in \{0,1\}$
            \medskip
            \State \textbf{Code for every process:}
            \State send $\langle input, b\rangle$ to the third party 
            \State \textbf{upon} receiving $\langle output, (b,g)\rangle$ from third party 
            \State \quad return $(b,g)$
            \medskip
            \State \textbf{Code for third party:}
            \State $bind,grade,call \gets \bot,\bot,(\bot)^{Id}$
            \State \textbf{upon} receiving $\langle input, b\rangle$ from a process $id$:
            \State \quad $call(id) \gets b$
            \State \textbf{upon} $|F| + \{id \in Id \setminus F \mid call(id) = b\} \ge n-f$:
            \State \quad $bind \gets b$
            \State \quad \textbf{if} $\{id \in Id \setminus F \mid call(id) = 1-b\} = 0$:
            \State \quad\quad $grade \gets A$
            \State \quad \textbf{else}
            \State \quad\quad $grade \gets C$
            \State \textbf{upon} $bind = \bot$ and $grade = A$ and not sent $\langle output, \cdot \rangle$ to process $id$:
            \State \quad send $\langle bind, B\rangle$ or $\langle output, (bind, A)\rangle$ to process $id$
            \State \textbf{upon} $bind = \bot$ and $grade = C$ and not sent $\langle output, \cdot \rangle$ to process $id$:
            \State \quad send $\langle bind, B\rangle$ or $\langle output, (\bot, C)\rangle$ to process $id$
        \end{algorithmic}
    \end{algorithm}
\end{protocol}

\begin{protocol}\label{prot:ImplGBCA}
    \lean{GBCA.Impl}
    %\leanok
    Parameters: $n$ processes with at most $f < n/3$ corrupted. Each process $p_i$ has input $b_i \in \{0,1\}$.
    \begin{algorithm}
        \caption{Implementation of GBCA taken from \cite{Abraham2022Efficient}}
        \begin{algorithmic}
            \State \textbf{Input:} $b \in \{0,1\}$
            \medskip
            \State send $\langle \textsc{input}, b_i \rangle$ to all processes
            \medskip
            \State \textbf{upon} receiving $\langle \textsc{input}, b \rangle$ from at least $f + 1$ processes and not sent $\langle \textsc{input}, b\rangle$ yet:
            \State \quad send $\langle \textsc{input}, b\rangle$ to all
            \medskip
            \State \textbf{upon} receiving $\langle \textsc{input}, b \rangle$ from at least $n - f$ processes:
            \State \quad $Approved \gets Approved \cup \{b\}$
            \State \quad \textbf{if} not sent $\langle \textsc{echo}, \cdot \rangle$ yet \textbf{then} send $\langle \textsc{echo}, b\rangle$ to all
            \medskip
            \State wait until 
            \State \quad (a) receiving $\langle \textsc{echo}, b \rangle$ from at least $n - f$ processes \textbf{or}
            \State \quad (b) receiving $\langle \textsc{echo}, \cdot \rangle$ from at least $n - f$ processes and $|Approved| > 1$
            \State \quad \textbf{if} (a) \textbf{then} send $\langle vote, b\rangle$ to all \textbf{else} send $\langle vote, \bot\rangle$ to all
            \medskip
            \State wait until 
            \State \quad (a) receiving $\langle \textsc{vote}, b \rangle$ from at least $n - f$ processes \textbf{or}
            \State \quad (b) receiving $\langle \textsc{vote}, \cdot \rangle$ from at least $n - f$ processes and $|Approved| > 1$
            \State \quad \textbf{if} (a) \textbf{then} send $\langle bind, b\rangle$ to all \textbf{else} send $\langle bind, \bot\rangle$ to all
            \medskip
            \State wait until 
            \State \quad (a) receiving $\langle \textsc{vote}, b \rangle$ from at least $n - f$ processes \textbf{or}
            \State \quad (b) receiving $\langle \textsc{vote}, \cdot \rangle$ from at least $n - f$ processes and there exists $b \in \{0,1\}$ such that $\langle bind, b\rangle$ has been received once and $\langle vote, b\rangle$ has been received from $t+1$ processes and $|Approved| > 1$ \textbf{or}
            \State \quad (c) receiving $\langle \textsc{vote}, \bot \rangle$ from at least $n - f$ processes and $|Approved| > 1$
            \State \quad \textbf{if} (a) \textbf{then} return $(b,A)$ \textbf{elif} (b) \textbf{then} return $(b,B)$ \textbf{else} return $(\bot, C)$
        \end{algorithmic}
    \end{algorithm}
\end{protocol}

\paragraph{Weak Common Coin (WCC)}
Processes start with no input and may return a bit.
A program implements WCC if it satisfies three properties: $\epsilon$-Goodness, Unpredictability, Termination.
\begin{itemize}
    \item $\epsilon$-Goodness: with probability $\epsilon$, all correct processes return 0 (resp. 1).
    \item Unpredictability: Until $t+1$ processes have called WCC, all outcomes are still possible.
    \item Termination: if $n-f$ correct processes take part to the protocol then all correct processes will eventually return.
\end{itemize}

\begin{protocol}\label{prot:SpecWCC}
    \lean{WCC.Spec}
    %\leanok
    Blablabla
    \begin{algorithm}
        \caption{Ideal Implementation of WCC}
        \begin{algorithmic}
            \State \textbf{Input:} none
            \medskip
            \State \textbf{Code for every process:}
            \State send $\langle input, b\rangle$ to the third party 
            \State \textbf{upon} receiving $\langle output, b\rangle$ from third party 
            \State \quad return $b$
            \medskip
            \State \textbf{Code for third party:}
            \State $r \gets \mu$ \Comment{$\mu: 1 \mapsto \epsilon; 0 \mapsto \epsilon; \bot \mapsto 1-2\epsilon$}
            \State \textbf{if} $r \ne \bot$:
            \State \quad send $\langle output, r \rangle$ to all
            \State \textbf{alse}
            \State \quad for all $id \in Id$, send $\langle output, 0\rangle$ \textbf{or} $\langle output, 1\rangle$ to process $id$
        \end{algorithmic}
    \end{algorithm}
\end{protocol}

\begin{protocol}\label{prot:ImplWCC}
    \lean{WCC.Impl}
    %\leanok
    Requires Gather and RSD
    \begin{algorithm}
        \caption{Implementation of WCC (yet to be determined)}
        \begin{algorithmic}
            \State $x \gets 0$
            \If{$y > 0$}
                \State loop
            \EndIf
        \end{algorithmic}
    \end{algorithm}
\end{protocol}

\begin{definition}\label{def:WCCConcrete}
    \lean{WCC.Concrete}
    %\leanok
    \uses{prot:ImplWCC, def:GatherConcrete, def:RSDConcrete}
    $\mathsf{WCC.Concrete} := \mathsf{WCC.Impl}[\mathsf{Gather.Concrete}, \mathsf{RSD.Concrete}]$, the instantiation of $\mathsf{WCC.Impl}$ with fully concrete sub-protocols.
\end{definition}

\begin{definition}\label{def:WCCHybrid}
    \lean{WCC.Hybrid}
    %\leanok
    \uses{prot:ImplWCC, prot:SpecGather, prot:SpecRSD}
    $\mathsf{WCC.Hybrid} := \mathsf{WCC.Impl}[\mathsf{Gather.Spec}, \mathsf{RSD.Spec}]$, the instantiation of $\mathsf{WCC.Impl}$ with ideal specifications of the sub-protocols.
\end{definition}

\subsection{WCC Sub-Protocols}

\paragraph{Gather} (taken from \cite{Stern2021Living,SterA2024Gather,Attiya2025Beyond})
Processes start with an input $x \in X$ and may return a partial function $f: Proc \rightharpoonup X$.
A program implements Gather if it satisfies four properties: Validity, Agreement, Common Core, Binding, Termination.
\begin{itemize}
    \item Validity: if a correct process returns $f$ and $f(q)$ is defined and $q$ is correct then $f(q)$ is $q$'s input.
    \item Agreement: if two correct processes return $f$ and $g$ then $f(q)$ and $g(q)$ are equal or one of them is undefined.
    \item Binding Common Core: Once a correct process has returned, there exists $S \subseteq Proc$ of size at least $n-f$ such that all $f$ returned by a correct process in the future is defined on $S$.
    \item Termination: if $n-f$ correct processes take part to Gather then all correct processes will eventually return.
\end{itemize}

\begin{protocol}\label{prot:SpecGather}
    \lean{Gather.Spec}
    %\leanok
    Blablabla
    \begin{algorithm}
        \caption{Ideal Implementation of Gather}
        \begin{algorithmic}
            \State \textbf{Input:} $m$
            \medskip
            \State \textbf{Code for every process:}
            \State send $\langle input, m\rangle$ to the third party 
            \State \textbf{upon} receiving $\langle output, f\rangle$ from third party:
            \State \quad return $f$
            \medskip
            \State \textbf{Code for third party:}
            \State $call, bind \gets \emptyset, \emptyset$
            \State \textbf{upon} receiving $\langle input, m\rangle$ from process $id$: 
            \State \quad $call \gets call \cup \{(id,m)\}$
            \State \textbf{if} $|F \cup \{id \in Id \setminus F \mid call[id] \ne \bot\}| > n-f$ and $bind = \bot$:
            \State \quad $bind \gets S \subseteq F \cup \{id \in Id \setminus F \mid (id,\cdot) \in call\}$ \Comment{Non-deterministic affectation}
            \State \textbf{upon} $bind \ne \bot$:
            \State \quad for all $id \in Id$, 
                            select $bind \subseteq S \subseteq F \cup \{id \in Id \setminus F \mid (id,\cdot) \in call\}$ and \\
                            send $\langle output, \{(id,m) \in call \mid id \in S \setminus F\} \cup \{(id, \cdot) \mid id \in S \cap F\}\rangle$ to process $id$
        \end{algorithmic}
    \end{algorithm}
\end{protocol}

\begin{protocol}\label{prot:ImplGather}
    \lean{Gather.Impl}
    %\leanok
    Might require BRB 
    \begin{algorithm}
        \caption{Implementation of Gather from \cite{Attiya2025Beyond}}
        \begin{algorithmic}
            \State \textbf{Input:} $b \in \{0,1\}$
            \State \textbf{Sub-Protocols:} $|Id|$ instances of BRB (process $id$ is leader in $BRB_{id}$)
            \medskip
            \State $BRB_{id}.call(m)$
            \medskip
            \State \textbf{upon} $BRB_i.return(m')$:
            \State \quad $AP \gets AP \cup \{(i,m')\}$
            \medskip
            \State \textbf{upon} $|AP| > n-f$ and not sent $\langle \textsc{echo}, \cdot\rangle$:
            \State \quad send $\langle echo, AP\rangle$ to all
            \medskip
            \State \textbf{upon} receiving approved $\langle \textsc{echo}, AP_{id} \rangle$ from at least $n - f$ processes and not sent $\langle \textsc{vote}, \cdot\rangle$:
            \State \quad send $\langle \textsc{vote}, \bigcup AP_{id} \rangle$
            \medskip
            \State \textbf{upon} receiving approved $\langle \textsc{vote}, AP_{id} \rangle$ from at least $n - f$ processes and not sent $\langle \textsc{bind}, \cdot\rangle$:
            \State \quad send $\langle \textsc{bind}, \bigcup AP_{id} \rangle$
            \medskip
            \State \textbf{upon} receiving approved $\langle \textsc{bind}, AP_{id} \rangle$ from at least $n - f$ processes:
            \State \quad return $\bigcup AP_{id}$
        \end{algorithmic}
    \end{algorithm}
\end{protocol}

\begin{definition}\label{def:GatherConcrete}
    \lean{Gather.Concrete}
    %\leanok
    \uses{prot:ImplGather, prot:ImplBRB}
    $\mathsf{Gather.Concrete} := \mathsf{Gather.Impl}[\mathsf{BRB.Impl}]$, the instantiation of $\mathsf{Gather.Impl}$ with the concrete implementation of BRB.
\end{definition}

\begin{definition}\label{def:GatherHybrid}
    \lean{Gather.Hybrid}
    %\leanok
    \uses{prot:ImplGather, prot:SpecBRB}
    $\mathsf{Gather.Hybrid} := \mathsf{Gather.Impl}[\mathsf{BRB.Spec}]$, the instantiation of $\mathsf{Gather.Impl}$ with the ideal specification of BRB.
\end{definition}

\paragraph{Random Secret Draw (RSD)}
The RSD protocol is composed of two operations: Assignment and Retrieval. 
Assignment affects a random and secret value to each process and Retrieval provides these values to processes.
Moreover, the API of RSD provides a notification function for each process that records which processes have a value assigned.
\begin{itemize}
    \item Assignment Validity: if all correct processes call Assignment, then all correct processes will eventually be receive a notification for each correct process.
    \item Assignment Totality: if a correct process receives a notification for process $id$ then all correct processes will get such a notification.
    \item Randomness: If no correct process has called Assignment then the value returned by Retrieval for any process follows the uniform distribution.
    \item Secrecy: If no correct process has called Retrieval then the adversary has no information about the return value of Retrieval.
    \item Retrieval Termination: if all correct processes call Retrieval then the operation terminates for all correct processes 
\end{itemize}

\begin{protocol}\label{prot:SpecRSD}
    \lean{RSD.Spec}
    %\leanok
    Blablabla
    \begin{algorithm}
        \caption{Ideal Implementation of RSD}
        \begin{algorithmic}
            \State $x \gets 0$
            \If{$y > 0$}
                \State loop
            \EndIf
        \end{algorithmic}
    \end{algorithm}
\end{protocol}

\begin{protocol}\label{prot:ImplRSD}
    \lean{RSD.Impl}
    %\leanok
    Require AVSS and BRB
    \begin{algorithm}
        \caption{Implementation of RSD from \cite{Freitas2022Distributed}}
        \begin{algorithmic}
            \State \textbf{Input:} $b \in \{0,1\}$
            \State \textbf{Sub-Protocols:} $|Id|$ instances of BRB (process $id$ is leader in $BRB_{id}$) and $|Id|^2$ instances of AVSS (process $id$ is dealer in $AVSS_{id}^{k}$)
            \medskip
            \State \textbf{Operation:} $RSD.Sh()$
            \State $retrieveEnabled, S \gets \bot, \emptyset$
            \State for all $id' \in Id$, $x_{id'} \gets \mu$ \Comment{$\mu$: Uniform distribution over the domain of secerts}
            \State $AVSS^{id'}_{id}.Sh.call(x_{id'})$
            \State \textbf{upon} $AVSS_{id'}^{id}.Sh.return()$:
            \State \quad $S \gets S \cup \{id'\}$
            \State \textbf{upon} $|AP| > n-f$ and not called $BRB_{id}$:
            \State \quad $BRB_{id}.call(S)$
            \State \textbf{upon} $BRB_{id'}.return(S')$ and $\forall k \in S', AVSS^{id'}_k.Sh.return()$:
            \State \quad $src[id'] \gets S'$
            \State \quad \textbf{event} $ValueAssigned(id')$
            \State \quad %\textbf{if} $retrieveEnabled$ \textbf{then} 
            $\forall k \in S', AVSS^{id'}_k.Rec.call()$
            \medskip
            %\State \textbf{Operation:} $RSD.EnableRec()$
            %\State $retrieveEnabled \gets \top$
            %\State $\forall id' \in Id$ \textbf{if} $ValueAssigned(id')$ \textbf{then} $AVSS.Rec.call()$
            %\medskip
            \State \textbf{Operation:} $RSD.Rec()$
            \State return $id' \mapsto \sum_{k \in src[id']}AVSS_k^{id'}.Rec.return()$
        \end{algorithmic}
    \end{algorithm}
\end{protocol}

\begin{definition}\label{def:RSDConcrete}
    \lean{RSD.Concrete}
    %\leanok
    \uses{prot:ImplRSD, prot:ImplBRB, prot:SpecAVSS}
    $\mathsf{RSD.Concrete} := \mathsf{RSD.Impl}[\mathsf{BRB.Impl}, \mathsf{AVSS.Spec}]$, the instantiation of $\mathsf{RSD.Impl}$ with the concrete implementation of BRB. $\mathsf{AVSS}$ is used as a black-box, so its specification is plugged in directly.
\end{definition}

\begin{definition}\label{def:RSDHybrid}
    \lean{RSD.Hybrid}
    %\leanok
    \uses{prot:ImplRSD, prot:SpecBRB, prot:SpecAVSS}
    $\mathsf{RSD.Hybrid} := \mathsf{RSD.Impl}[\mathsf{BRB.Spec}, \mathsf{AVSS.Spec}]$, the instantiation of $\mathsf{RSD.Impl}$ with the ideal specifications of BRB and AVSS.
\end{definition}

\subsection{Basic Primitives}

\paragraph{Byzantine Reliable Broadcast (BRB)} (taken from \cite{Bracha1987Asynchronous})
Each instance of BRB has a designated \emph{leader} with an input message $m$ and processes may return a received message.
A program implements BRB if it satisfies three properties: Validity, Totality and Agreement.
\begin{itemize}
    \item Validity: if the leader is not corrupted then all correct processes must return its input.
    \item Agreement: if two correct processes return then they return the same value.
    \item Totality: if a correct process returns then all correct processes eventually return.
\end{itemize}

\begin{protocol}\label{prot:SpecBRB}
    \lean{BRB.Spec}
    %\leanok
    Blablabla
    \begin{algorithm}
        \caption{Ideal Implementation of BRB}
        \begin{algorithmic}
            \State \textbf{Input:} $m \in \{0,1\}^*$ for the leader
            \medskip
            \State \textbf{Code for the leader:}
            \State send $\langle input, m\rangle$ to the third party 
            \medskip
            \State \textbf{Code for every process:}
            \State \textbf{upon} receiving $\langle output, m\rangle$ from third party:
            \State \quad return $m$
            \medskip
            \State \textbf{Code for third party:}
            \State \textbf{upon} receiving $\langle input, m\rangle$ from leader:
            \State \quad \textbf{if} leader is correct:
            \State \quad\quad send $\langle output, m\rangle$ to all
            \State \quad \textbf{else}
            \State \quad\quad send $\langle output, m'\rangle$ to all \Comment{Arbitrary $m'$}
        \end{algorithmic}
    \end{algorithm}
\end{protocol}

\begin{protocol}\label{prot:ImplBRB}
    \lean{BRB.Impl}
    %\leanok
    Requires no sub-protocol. Parameters: $n$ processes with at most $f < n/3$ corrupted; leader $\ell$ with input $m$.
    \begin{algorithm}
        \caption{Implementation of BRB taken from \cite{Bracha1987Asynchronous}}
        \begin{algorithmic}
            \State \textbf{Code for the leader $\ell$ with input $m$:}
            \State \quad send $\langle \textsc{init}, m \rangle$ to all processes
            \medskip
            \State \textbf{Code for every process $p_i$:}
            \State \textbf{upon} first reception of $\langle \textsc{init}, m \rangle$ from $\ell$ \textbf{do}
                \State \quad send $\langle \textsc{echo}, m \rangle$ to all processes
            \State \textbf{upon} reception of $\langle \textsc{echo}, m \rangle$ from at least $n-f$ distinct processes \textbf{do}
                \State \quad \textbf{if} $\langle \textsc{ready}, \cdot \rangle$ has not yet been sent \textbf{then}
                    \State \quad\quad send $\langle \textsc{ready}, m \rangle$ to all processes
            \State \textbf{upon} reception of $\langle \textsc{ready}, m \rangle$ from at least $f+1$ distinct processes \textbf{do}
                \State \quad \textbf{if} $\langle \textsc{ready}, \cdot \rangle$ has not yet been sent \textbf{then}
                    \State \quad\quad send $\langle \textsc{ready}, m \rangle$ to all processes
            \State \textbf{upon} reception of $\langle \textsc{ready}, m \rangle$ from at least $2f+1$ distinct processes \textbf{do}
                \State \quad deliver $m$
        \end{algorithmic}
    \end{algorithm}
\end{protocol}


\paragraph{Asynchronous Verifiable Secret Sharing (AVSS)}
AVSS is actually a pair of protocol $(Sh, Ret)$ with a designated \emph{dealer} owning a secret $s$.
A program implements a $(1-\epsilon)$-correct AVSS if it satisfies six properties: Valid Termination, Total Termination, $Rec$ Termination, Binding, Secrecy.
\begin{itemize}
    \item Validity: if the dealer is correct and a correct process returns from $Ret$ then the return value is the dealer's secret.
    \item Sharing Termination: if the dealer is correct and calls $Sh$ then all correct processes will enventually return from $Sh$.
    \item Sharing Totality: if a correct process returns from $Sh$ then all correct processes will eventually return from $Sh$.
    \item Retrieval Termination: if all correct processes call $Ret$ then all correct processes will eventually return from $Ret$.
    \item Commitment: if a correct process returns from $Sh$ then there exists a secret $s$ such that correct process can only retrieve $s$ from that point.
    \item Secrecy: if the dealer is correct and no correct process has called $Ret$ then the adversary has no information about the future retrived value.
\end{itemize}

\begin{protocol}\label{prot:SpecAVSS}
    \lean{AVSS.Spec}
    %\leanok
    Blablabla
    \begin{algorithm}
        \caption{AVSS Specification}
        \begin{algorithmic}
            \State $x \gets 0$
            \If{$y > 0$}
                \State loop
            \EndIf
        \end{algorithmic}
    \end{algorithm}
\end{protocol}

AVSS is used as a blackbox so no implementation.

\subsection{Summary of Properties}

The linear safety properties are preserved by Segala's probabilistic forward simulation. 
The branching safety properties can be transformed into linear ones by enhancing the 
interface with an external binding label.
Segala's simulation does not preserve liveness by default. 
However, it can be strengthen following \cite{Dongol2022Weak} to preserve liveness 
in a complete manner.

\begin{table}[ht]
    \centering
    \begin{tabular}{|c||c|c|c|c|c|c|}
        \hline
        & Linear Safety & Branching Safety & Liveness \\
        \hline
        ABA & Validity \& Agreement & &  \\
        GBCA & Validity \& Agreement & Binding & Termination \\
        WCC & & & Termination  \\
        Gather & Validity \& Agreement & Binding Common Core & Termination  \\
        RSD & & & \makecell{Assignment Validity \& \\ Assignment Totality \& \\ Retrieval Termination} \\
        BRB & Validity \& Agreement & & Totality \\
        AVSS & Validity & Commitment & \makecell{Sharing Termination \& \\ Sharing Totality \& \\ Retrieval Termination} \\
        \hline
    \end{tabular}
    \caption{Table of CTL Properties}
\end{table}

Segala's simulation was designed to be sound and complete for trace distribution so
it preserves probabilistic safety propreties.
Note that Unpredictability is a probabilistic safety property.
For probabilistic liveness, we will rely on the liveness preservation as it is sound 
(but not complete). 
The lack of completeness is not an isue because the limitations of the termination
proof rule is probabilistic convergence which is never dealt with by our simulations.
The probabilistic convergence argument will only be necessary to prove that the ideal 
implementation of ABA almost-surely terminates and can be discharged to the 
super-Martingale proof rule of \cite{Enea2026Verifying}.
For Secrecy, we use the same trick as for branching: enhancing the interface. 


\begin{table}[ht]
    \centering
    \begin{tabular}{|c||c|c|c|c|c|c|}
        \hline
        & Probabilistic Safety & Probabilistic Liveness & Secrecy \\
        \hline
        ABA & & Almost-Sure Termination & \\
        GBCA & & & \\
        WCC & $\epsilon$-Goodness \& Unpredictability & & \\
        Gather & & & \\
        RSD & Randomness & & Secrecy \\
        BRB & & & \\
        AVSS & & & Secrecy \\
        \hline
    \end{tabular}
    \caption{Table of non-CTL Properties}
\end{table}



